% Irodov VI §6.3 originals 135--147, printed p248, PDF p257.
% Answers user PDF pages 354--355.
\atomictermsproblem{135}{248}
两种元素的特征$K_\alpha$线波长分别为
$250\,\mathrm{pm}$和$179\,\mathrm{pm}$。
求它们在周期表中间隔着多少个元素。
\exsolution
莫塞莱定律给
$\lambda_{K\alpha}=hc/[(3/4)\mathrm{Ry}(Z-1)^2]$，
所以
\[
 Z=1+\sqrt{\frac{hc}{(3/4)\mathrm{Ry}\lambda}}.
\]
对$250\,\mathrm{pm}$得$Z_1\simeq23$（钒），
对$179\,\mathrm{pm}$得$Z_2\simeq27$（钴）。
中间的元素序数为24、25、26，故
\[
 \boxed{N=Z_2-Z_1-1=3}.
\]
\excheck “位于两者之间”不包含端点元素，
因而是三种而不是四种。

\atomictermsproblem{136}{248}
以镍（$Z=28$）为靶的X射线管，
其$K_\alpha$谱线与连续谱短波限的波长差
为$\Delta\lambda=84\,\mathrm{pm}$。
求加速电压$U$。
\exsolution
莫塞莱近似给
$K_\alpha$光子能量
$\Delta E_\alpha=(3/4)\mathrm{Ry}(27)^2$，
故其波长$\lambda_\alpha=hc/\Delta E_\alpha
\simeq167\,\mathrm{pm}$。
连续谱短波限更短，
$\lambda_{\min}=\lambda_\alpha-\Delta\lambda
\simeq83\,\mathrm{pm}$。
由$eU=hc/\lambda_{\min}$，
\[
 \boxed{U=\frac{hc}{e(\lambda_\alpha-\Delta\lambda)}
 \simeq15\,\mathrm{kV}}.
\]
\excheck 若误把短波限设为$\lambda_\alpha+\Delta\lambda$，
会得到不合理的低电压；短波限应位于更短波一侧。

\atomictermsproblem{137}{248}
铝靶X射线管的连续谱短波限
为$\lambda_{\min}=0.50\,\mathrm{nm}$，
而其K系激发电压为$U_K=1.56\,\mathrm{kV}$。
问这时还能观察到特征K系谱线吗？
\exsolution
连续谱短波限对应最高电子动能
$eU=hc/\lambda_{\min}$。
因此
\[
 U=\frac{hc}{e(0.50\,\mathrm{nm})}
 \simeq2.48\,\mathrm{kV}.
\]
因$2.48\,\mathrm{kV}>1.56\,\mathrm{kV}$，
入射电子足以从K层击出电子，
随后更高层电子填空位便发出K系特征线，
\[
 \boxed{\text{可以观察到K系谱线。}}
\]
\excheck 产生连续谱不自动保证能产生某一特征系，
要比较实际电子能量与对应内层电离阈值。

\atomictermsproblem{138}{248}
某X射线管的电压从$U_1=10\,\mathrm{kV}$
升高为$U_2=20\,\mathrm{kV}$。
底本将所比较的固定特征波长称为$K_\infty$线，
并说它与连续谱短波限的波长间距
增为原来的$\eta=3.0$倍。
求靶元素原子序数$Z$，并核查底本答案。
\exsolution
先单独解决\emph{几何上的波长间距}问题。
两电压下连续谱短波限
\[
 \lambda_i=\frac{hc}{eU_i},\qquad
 \lambda_1=123.984\,\mathrm{pm},\quad
 \lambda_2=61.992\,\mathrm{pm}.
\]
令待定特征波长为$\lambda_X$。
若初始$10$kV时该特征线已经能被激发，
必须有$\lambda_X\ge\lambda_1$，
所以波长间隔满足
\[
 \lambda_X-\lambda_2
 =\eta(\lambda_X-\lambda_1),\qquad
 \boxed{\lambda_X=
 \frac{\eta\lambda_1-\lambda_2}{\eta-1}
 =154.980\,\mathrm{pm}}.
\]
它对应的光子能量为
$\varepsilon_X=hc/\lambda_X=8.00\,\mathrm{keV}$。
\emph{关键是辨认特征线的种类：}

若按底本书末采用的\emph{莫塞莱$K_\alpha$线}
近似能量
$\varepsilon_{K_\alpha}
=\frac34\mathrm{Ry}(Z-1)^2$，则
\[
 \boxed{Z=1+\sqrt{\frac{4\varepsilon_X}{3\mathrm{Ry}}}
 \simeq29.0},\qquad
 \boxed{\text{靶材为铜（Cu）}}.
\]
这与书末$Z=29$及其所列含$3\mathrm{Ry}/4$
的公式相符。
但原题\emph{字面写$K_\infty$}，
它是K系极限（对应K层电离边），
最简氢样屏蔽模型应使用
$\varepsilon_{K_\infty}\simeq\mathrm{Ry}(Z-1)^2$，
因而
\[
 \boxed{Z_{K_\infty}\simeq
 1+\sqrt{\frac{8.00\times10^3}{13.606}}
 \simeq25.25\approx25}
\]
（模型精度内对应锰），而非29。
再者，若把波长“间距”理解为绝对值且
\emph{没有}要求第一电压下产生特征线，
还存在跨越两短波限的代数分支：
\[
 |\lambda_X-\lambda_2|=3|\lambda_X-\lambda_1|,
 \quad
 \lambda_2<\lambda_X<\lambda_1
 \ \Longrightarrow\
 \lambda_X=\frac{3\lambda_1+\lambda_2}{4}
 =108.486\,\mathrm{pm}.
\]
此时$\varepsilon_X=11.43$keV，
超过第一电压电子最高能量10keV，
第一状态不能激发该线。
\excheck 已对底本原题图像和书末第138题公式直接校核：
\emph{原题写$K_\infty$，书末解法却用了
$K_\alpha$的$\frac34$莫塞莱系数}。
所以保留\emph{书末口径}的$Z=29$，
同时列出严格按字面$K_\infty$解读的
$Z\approx25$与未激发分支；
不能仅用前版“书末答案错误”一句带过，
也不可把$K_\infty$静默改译成$K_\alpha$。

\atomictermsproblem{139}{248}
某金属X射线K、L两吸收限对应角频率之差为
$\Delta\omega=6.85\times10^{18}\,\mathrm{s^{-1}}$。
求该金属。
\exsolution
在莫塞莱模型中，
K、L结合能的差约为
$\Delta E=(3/4)\mathrm{Ry}(Z-1)^2$，
故
\[
 \Delta\omega=\frac{\Delta E}{\hbar}
 =\frac{3\mathrm{Ry}}{4\hbar}(Z-1)^2.
\]
反解
\[
 \boxed{Z=1+\sqrt{\frac{4\hbar\Delta\omega}
 {3\mathrm{Ry}}}\simeq22}.
\]
原子序数22是$\boxed{\mathrm{Ti}\text{（钛）}}$。
\excheck 题目给的是\emph{角频率}，
因此用$\Delta E=\hbar\Delta\omega$，不是$h\Delta\omega$。

\atomictermsproblem{140}{248}
钒（$Z=23$）的L层吸收限波长
$\lambda_L=2.4\,\mathrm{nm}$。
求其K层电子结合能。
\exsolution
L层结合能为
$E_L=hc/\lambda_L\simeq0.517\,\mathrm{keV}$。
莫塞莱近似给K、L层结合能差
$\Delta E_{KL}=(3/4)\mathrm{Ry}(Z-1)^2$，
对钒$Z-1=22$，
$\Delta E_{KL}\simeq4.94\,\mathrm{keV}$。
所以
\[
 \boxed{E_K=E_L+\Delta E_{KL}\simeq5.46\,\mathrm{keV}
 \approx5.5\,\mathrm{keV}}.
\]
\excheck 不能直接用K$\alpha$能量当K电子全部结合能，
因为该谱线对应两层能量\emph{之差}。

\atomictermsproblem{141}{248}
已知钛的K系第一条谱线$K_\alpha$
与本系短波极限的波长差为
$\Delta\lambda=26\,\mathrm{pm}$。
求钛L层电子结合能。
\exsolution
钛$Z=22$。
莫塞莱关系给K$\alpha$能量
$E_\alpha=(3/4)\mathrm{Ry}(21)^2
\simeq4.50\,\mathrm{keV}$，
所以$\lambda_\alpha=hc/E_\alpha\simeq275.5\,\mathrm{pm}$。
K系短波极限对应K层结合能，
其波长
$\lambda_K=\lambda_\alpha-\Delta\lambda
\simeq249.5\,\mathrm{pm}$。
于是
$E_K=hc/\lambda_K\simeq4.97\,\mathrm{keV}$，
$E_L=E_K-E_\alpha$，故
\[
 \boxed{E_L\simeq0.47\,\mathrm{keV}
 \approx0.5\,\mathrm{keV}}.
\]
\excheck 两谱线间隔是波长差，不等于能量差；
计算必须分别取$hc/\lambda$再相减。

\atomictermsproblem{142}{248}
锌的$K_\alpha$辐射照射铁，
从铁K层打出电子。
已知铁的K吸收限波长$\lambda_K=174\,\mathrm{pm}$。
求光电子动能和速率。
\exsolution
锌$Z=30$的K$\alpha$光子能量
$\varepsilon=(3/4)\mathrm{Ry}(29)^2
\simeq8.58\,\mathrm{keV}$。
铁K层结合能
$E_{b,\mathrm{Fe}}=hc/(174\,\mathrm{pm})
\simeq7.13\,\mathrm{keV}$。
光电子动能
\[
 \boxed{T=\varepsilon-E_b\simeq1.45\,\mathrm{keV}}.
\]
在此能区可用非相对论近似
$v=\sqrt{2T/m_e}$，故
\[
 \boxed{v\simeq2.26\times10^7\,\mathrm{m/s}}.
\]
\excheck 必须用锌的特征谱线能量减去\emph{铁}的K层结合能。

\atomictermsproblem{143}{248}
求朗德$g$因子：(a) S项；(b)单重态。
\exsolution
一般LS耦合有
\[
 g_J=1+\frac{J(J+1)+S(S+1)-L(L+1)}
 {2J(J+1)}.
\]
(a) S项$L=0,J=S>0$，
故$\boxed{g_J=2}$。
若同时$S=0$则$J=0$，上述公式为$0/0$，
磁矩本身为零。
(b) 单重态$S=0$、$J=L>0$，
故$\boxed{g_J=1}$。
\excheck 单重S态的$J=0$使$g$的常规定义不适用，
不能强行把$0/0$写成某个数。

\atomictermsproblem{144}{248}
求下列光谱项的朗德因子：
(a)\({}^{6}F_{1/2}\)；
(b)\({}^{4}D_{1/2}\)；
(c)\({}^{5}F_2\)；
(d)\({}^{5}P_1\)；
(e)\({}^{3}P_0\)。
\exsolution
代入
$g=1+[J(J+1)+S(S+1)-L(L+1)]/[2J(J+1)]$。
(a)$(S,L,J)=(5/2,3,1/2)$，故$g=-2/3$；
(b)$(3/2,2,1/2)$，得$g=0$；
(c)$(2,3,2)$，得$g=1$；
(d)$(2,1,1)$，得$g=5/2$；
(e)$J=0$，公式分母为零、磁矩为零，
$g$不作通常定义。于是
\[
 \boxed{g_a=-\tfrac23,\quad g_b=0,\quad
 g_c=1,\quad g_d=\tfrac52,\quad g_e:\text{未定义}}.
\]
\excheck $g$可为负，零总角动量态不应硬算$g$的数值。

\atomictermsproblem{145}{248}
求以下原子磁矩模，以玻尔磁子为单位：
(a)单重F态；
(b)\({}^{2}D_{3/2}\)；
(c)$S=1,L=2,g_J=4/3$的状态。
\exsolution
取
$\mu=g_J\mu_B\sqrt{J(J+1)}$。
(a) $S=0,L=J=3,g=1$，
\[
 \boxed{\mu_a=\sqrt{12}\,\mu_B=2\sqrt3\,\mu_B}.
\]
(b) $S=1/2,L=2,J=3/2$，
$g=1+[(15/4)+(3/4)-6]/[2(15/4)]=4/5$，故
\[
 \boxed{\mu_b=\frac45\sqrt{\frac{15}4}\,\mu_B
 =\frac{2\sqrt{15}}5\,\mu_B}.
\]
(c) $S=1,L=2$，由$g_J=4/3$反求$J=3$，
因此
\[
 \boxed{\mu_c=\frac43\sqrt{12}\,\mu_B
 =\frac{8\sqrt3}{3}\,\mu_B}.
\]
\excheck (c)必须先从$g$确定$J$，
不能任取$J=1$或2。

\atomictermsproblem{146}{248}
一个原子的谱项为$D_2$，
磁矩沿外场的最大投影为$4\mu_B$。
求其总自旋角动量模。
\exsolution
$D$给$L=2$，下标给$J=2$。
最大磁矩投影$g_JJ\mu_B=4\mu_B$，
故$g_J=2$。
朗德公式化简为
\[
 g_J=1+\frac{S(S+1)}{2J(J+1)}
 =1+\frac{S(S+1)}{12}.
\]
令$g_J=2$得$S(S+1)=12$，
即$S=3$，所以
\[
 \boxed{M_S=\hbar\sqrt{S(S+1)}=
 2\sqrt3\,\hbar}.
\]
\excheck 此题使用的是磁矩\emph{沿场最大投影}，
不是磁矩模$g\sqrt{J(J+1)}\mu_B$。

\atomictermsproblem{147}{248}
某原子$L=2,S=1$，处于弱磁场中。
已知总角动量与磁场方向的最小夹角为$30^\circ$。
求原子磁矩的模。
\exsolution
空间量子化下总角动量最大投影
$M_{J,z}=J\hbar$，而模为
$M_J=\hbar\sqrt{J(J+1)}$，
故最小夹角满足
\[
 \cos\theta_{\min}=
 \frac J{\sqrt{J(J+1)}}=
 \sqrt{\frac J{J+1}}.
\]
取$\theta_{\min}=30^\circ$，
解得$J=3$。
对$L=2,S=1,J=3$，
$g_J=4/3$，因此
\[
 \boxed{\mu=\frac43\mu_B\sqrt{3\cdot4}
 =\frac{8\sqrt3}{3}\,\mu_B}.
\]
\excheck 磁矩模与沿磁场投影不同；
前者有$\sqrt{J(J+1)}$系数。
